\chapter{Introduction}

\section{Overview and Background}
Active recall, the practice of retrieving information from memory instead of simply re-reading it, is one of the most effective study techniques available to students. Quizzes are its most natural form. In practice, however, learners rarely have a quiz that matches exactly what they are studying: published question banks are generic, and writing good questions manually, with plausible wrong options, model answers and test cases, takes longer than studying the material itself.

Recent large language models (LLMs) can read a passage and produce well-formed questions about it. The remaining engineering problem is to make this reliable and useful: the model must be given trustworthy source material, its output must be forced into a machine-readable structure, the answers must be graded, and the whole experience must feel like a quiz and not a chat window.

\textbf{Recall} is a web application that addresses this problem. A learner either types a topic or uploads study material, chooses a difficulty and a question mix, and receives a ready-to-take quiz containing multiple-choice, written and coding questions. The application runs on the learner's own machine, so uploaded notes never have to be stored on a third-party server other than the model provider that is explicitly selected.

\section{Problem Statement}
Students preparing for examinations face the following recurring difficulties:
\begin{itemize}
\item Ready-made question banks do not cover the learner's own notes, slides or syllabus.
\item Writing questions, distractors, model answers and test cases manually is time-consuming.
\item Asking a general chat assistant for ``a quiz'' produces free-form text that cannot be taken, scored, timed or reviewed inside an application.
\item Questions invented by a model without source material may contain factual errors, and the learner has no way to see where a question came from.
\item Written and coding answers need feedback, which most static quiz tools cannot give.
\end{itemize}
There is therefore a need for a single tool that generates \textit{grounded}, \textit{structured} and \textit{gradable} quizzes from arbitrary topics and documents.

\section{Objectives of the System}
The objectives of the project are:
\begin{enumerate}
\item To generate quizzes on any user-typed topic, grounded in web sources whose names and links are shown to the learner.
\item To generate quizzes from uploaded PDF, DOCX, PPTX, TXT and Markdown files and from pasted text.
\item To support three question types (multiple choice, written answer and coding) at Easy, Medium and Hard difficulty.
\item To enforce a strict JSON schema on model output and to recover automatically from malformed or empty replies.
\item To provide two engines, a fast cloud engine (Groq) and an offline local engine (Ollama), with automatic fallback.
\item To grade written answers with AI feedback and to execute coding answers against test cases.
\item To provide Practice mode (instant feedback) and Exam mode (timer, question grid, flag for review), together with history, statistics and export.
\end{enumerate}

\section{Scope and Limitations}
\textbf{Scope.} The system covers quiz generation, taking, grading, reviewing and exporting for a single learner using a modern desktop or mobile browser. It includes the Flask back end, the research and document-ingestion modules, the code sandbox and the single-page front end.

\textbf{Limitations.}
\begin{itemize}
\item The system depends on the quality of the selected language model; the small local model gives noticeably weaker questions than the cloud model.
\item Quiz history is stored only in the browser's local storage; there are no user accounts and no cross-device synchronisation.
\item The code runner executes Python only, and it runs submitted code in a plain subprocess on the host machine. It is intended for local, single-user use and is not a hardened multi-tenant sandbox.
\item Scanned PDFs without a text layer cannot be read, because no optical character recognition is performed.
\item The project does not currently include an automated test suite; testing was performed manually (Chapter~\ref{chap:security}).
\end{itemize}
